\documentclass[a4paper]{article}

% Basic packages
% Español (México) con XeLaTeX: fontspec para fuentes Unicode, polyglossia
% para la división silábica y los nombres del documento en la variante mexicana.
\usepackage{fontspec}
\usepackage{polyglossia}
\setdefaultlanguage[variant=mexican]{spanish}
% Los nombres chinos van como «pinyin (original)»: xeCJK da a los caracteres
% Han su propia fuente; sin ella XeLaTeX los omite con un aviso.
% FandolSong no cubre algunos caracteres de nombres (U+73FA); WenQuanYi Zen Hei sí.
\usepackage[AutoFallBack=true]{xeCJK}
\setCJKmainfont{FandolSong-Regular.otf}
\setCJKfallbackfamilyfont{\CJKrmdefault}{WenQuanYi Zen Hei}
% polyglossia no traduce los nombres de listings.
\AtBeginDocument{\providecommand{\lstlistingname}{}\renewcommand{\lstlistingname}{Listado}%
  \providecommand{\lstlistlistingname}{}\renewcommand{\lstlistlistingname}{Índice de listados}}
\usepackage{amsmath, amssymb}
\usepackage{graphicx}
\usepackage[margin=2.5cm]{geometry}
\usepackage[most]{tcolorbox}
\usepackage{etoolbox}
\usepackage{listings}
\usepackage{hyperref}
\usepackage{booktabs}
\usepackage{subcaption}
\usepackage{float}

% Highlight boxes
\newtcolorbox{knowledgebox}[1]{
    enhanced, colback=blue!5!white, colframe=blue!75!black, colbacktitle=blue!75!black,
    coltitle=white, fonttitle=\bfseries, title={#1}, attach boxed title to top left={yshift=-2mm, xshift=2mm},
    boxrule=1pt, sharp corners
}

\newtcolorbox{importantbox}[1]{
    enhanced, colback=yellow!10!white, colframe=yellow!80!black, colbacktitle=yellow!80!black,
    coltitle=black, fonttitle=\bfseries, title={#1}, sharp corners
}

\newtcolorbox{warningbox}[1]{
    enhanced, colback=red!5!white, colframe=red!75!black, colbacktitle=red!75!black,
    coltitle=white, fonttitle=\bfseries, title={#1}, sharp corners
}

% Listings
\lstset{
    language=Python,
    basicstyle=\ttfamily\small,
    keywordstyle=\color{blue},
    stringstyle=\color{red!60!black},
    commentstyle=\color{green!60!black},
    breaklines=true,
    frame=single,
    numbers=left,
    numberstyle=\tiny\color{gray},
    captionpos=b,
    extendedchars=true
}

% Front-page metadata
\newcommand{\notetitle}{Gemini Agentic Vision: Capacidades fundamentales del modelo y Post-Training}
\newcommand{\noteauthors}{Elaboradas a partir de materiales públicos del curso}
\newcommand{\notedate}{2025}
\newcommand{\videochannel}{Wudaokou Naxi (五道口纳什)}
\newcommand{\videopublishdate}{2025}
\newcommand{\videoduration}{aproximadamente 16 minutos}
\newcommand{\videourl}{}
\newcommand{\videocoverpath}{cover.jpg}
\newcommand{\repourl}{https://github.com/hqhq1025/ai-course-notes}


% Footer with repo info
\usepackage{fancyhdr}
\pagestyle{fancy}
\fancyhf{}
\fancyfoot[C]{\thepage}
\fancyfoot[R]{\footnotesize\color{gray}\href{\repourl}{AI Course Notes}}
\renewcommand{\headrulewidth}{0pt}
\renewcommand{\footrulewidth}{0pt}

\begin{document}

\begin{titlepage}
\centering
{\Large Notas del curso\par}
\vspace{1.2cm}
{\huge\bfseries \notetitle\par}
\vspace{0.8cm}
{\large \noteauthors\par}
\vspace{0.3cm}
{\large \notedate\par}
\vspace{1.2cm}

\ifdefempty{\videocoverpath}{
{\small Al generar las notas, indique la ruta local de la portada del video.\par}
}{
\includegraphics[width=0.82\textwidth,height=0.45\textheight,keepaspectratio]{\videocoverpath}\par
}

\vfill
\begin{tcolorbox}[width=0.9\textwidth, colback=black!2!white, colframe=black!60, sharp corners]
\textbf{Autor o canal del video}: \videochannel\par
\textbf{Fecha de publicación}: \videopublishdate\par
\textbf{Duración del video}: \videoduration\par
\ifdefempty{\videourl}{}{
\textbf{Enlace del video}: \href{\videourl}{\nolinkurl{\videourl}}\par
}
\par
\textbf{Página del proyecto}: \href{\repourl}{\nolinkurl{\repourl}}\par
\end{tcolorbox}
\end{titlepage}

\tableofcontents
\newpage

%% --- Inicio del contenido --- %%

\section{Introducción: el significado inspirador de Agentic Vision}

Esta entrega presenta la función Agentic Vision de Gemini 2.5 Flash, publicada recientemente por Google. Este trabajo combina el razonamiento visual con la ejecución de código, de modo que el modelo puede «investigar» activamente una imagen, en lugar de «echarle un vistazo y adivinar» de manera pasiva.

\begin{knowledgebox}{La relación entre Modern Agents y Agentic RL Training}
El expositor ubicó con claridad ambas series:
\begin{itemize}
    \item \textbf{Modern Agents}: se centra en el razonamiento, las llamadas a herramientas en múltiples turnos y las técnicas de Novel Prompting
    \item \textbf{Agentic RL Training}: se centra en los algoritmos, la arquitectura de sistemas y el Post-Training
\end{itemize}
Agentic Vision conecta precisamente ambas series: el efecto de razonamiento en múltiples turnos que muestra inspira la manera de reforzar e incentivar, mediante Post-Training, las capacidades básicas del modelo.
\end{knowledgebox}

\section{Mecanismos centrales de Agentic Vision}

\subsection{De la "observación estática" a la "investigación activa"}

Los VLM tradicionales procesan la entrada visual desde una \textbf{perspectiva estática única}: si se omite algún detalle sutil, solo queda adivinar. Agentic Vision transforma la comprensión de imágenes de una conducta estática a un \textbf{proceso agéntico}:

\begin{importantbox}{Las tres innovaciones de Agentic Vision}
\begin{itemize}
    \item \textbf{Razonamiento visual + ejecución de código}: el modelo escribe código Python para operar sobre la imagen
    \item \textbf{Tool Call de múltiples turnos}: el ciclo de ReAct de Think → Act → Observe
    \item \textbf{Borrador visual (Visual Scratchpad)}: el modelo genera de manera autónoma imágenes anotadas para apoyar su razonamiento
\end{itemize}
\end{importantbox}

\subsection{El único Tool: Code Execution}

El diseño del Tool en Agentic Vision es extremadamente sencillo: solo existe una \textbf{herramienta de ejecución de código} (Code Execution on Images). Al escribir código, el modelo puede lograr:

\begin{itemize}
    \item \textbf{Bounding Box}: dibujar cuadros de detección
    \item \textbf{Crop}: recortar regiones de la imagen
    \item \textbf{Rotate/Zoom}: rotar y ampliar
    \item \textbf{Annotate}: colocar anotaciones en la imagen (texto, flechas, numeración)
    \item \textbf{Draw}: dibujar líneas auxiliares
    \item \textbf{Save}: guardar la imagen procesada
\end{itemize}

\begin{warningbox}{Observación clave: el BBox no lo produce un modelo de CV externo}
En Agentic Vision, las coordenadas del Bounding Box las \textbf{produce el propio Gemini 2.5 Flash}: al escribir el código, pasa directamente las coordenadas del BBox. Se trata de una \textbf{capacidad visual básica nativa} del modelo (detección, segmentación, reconocimiento), y no de la dependencia de un modelo externo de detección de objetos.
\end{warningbox}

\subsection{El ciclo de ReAct: Think → Act → Observe}

El flujo de trabajo de Agentic Vision es un proceso estándar de ReAct (Reasoning + Acting):

\begin{enumerate}
    \item \textbf{Think}: analiza la pregunta del usuario y planea los pasos de operación visual
    \item \textbf{Act}: genera código Python y llama al Code Execution Tool
    \item \textbf{Observe}: revisa el resultado de la ejecución del código (la imagen procesada) y decide si se necesita otra operación
    \item Se repite el ciclo hasta llegar a la respuesta final
\end{enumerate}

\begin{knowledgebox}{Esto surge del Post-Training}
El ciclo de ReAct no está codificado de manera rígida mediante un marco externo, sino que surge de forma espontánea en el modelo, mediante Post-Training (aprendizaje por refuerzo), como una conducta de llamadas a herramientas de múltiples turnos. El modelo aprende en qué momento debe hacer zoom in, recortar o anotar.
\end{knowledgebox}

\section{Comparación con el Agente SAM3}

\subsection{Construcción explícita vs. surgimiento autónomo}

El Agente SAM3 presentado en la entrega anterior también implementó Visual Prompting, pero de una manera completamente distinta:

\begin{table}[H]
\centering
\begin{tabular}{p{3cm}p{5cm}p{5cm}}
\toprule
\textbf{Dimensión} & \textbf{SAM3 + VLM} & \textbf{Gemini Agentic Vision} \\
\midrule
Herramienta visual & SAM3 (modelo de segmentación especializado) & Code Execution (de propósito general) \\
Arquitectura & Agent Pipeline construido explícitamente & Conducta incentivada por Post-Training \\
Boceto visual & Generado por SAM3 mediante una superposición de Mask & El modelo escribe el código de manera autónoma \\
Origen del BBox & Modelo/herramienta externos & Capacidad nativa del modelo \\
Flexibilidad & Limitada por las Tool definidas & En teoría, puede realizar cualquier operación sobre la imagen \\
\bottomrule
\end{tabular}
\caption{Comparación entre el Agente SAM3 y Gemini Agentic Vision}
\end{table}

\subsection{La evolución del Visual Prompting}

Ambos esquemas producen un «boceto visual» (Visual Scratchpad), pero de origen distinto:

\begin{itemize}
    \item Agente SAM3: los humanos \textbf{diseñaron explícitamente} el proceso de segmentar, anotar y devolver el resultado al VLM
    \item Agentic Vision: el modelo \textbf{decide de manera autónoma} qué tipo de boceto visual necesita para apoyar su razonamiento
\end{itemize}

\subsection{Resumen de la sección}

De la construcción explícita al surgimiento autónomo, Agentic Vision representa la dirección de evolución de la técnica de Visual Prompting: dejar que el modelo decida por sí mismo cómo "ver" la imagen, en lugar de que los humanos predefinan la manera de observarla.

\section{Relación entre las capacidades fundamentales y el post-entrenamiento}

\subsection{El paradigma de dos etapas}

Agentic Vision revela un paradigma de entrenamiento importante:

\begin{importantbox}{Capacidades fundamentales + liberación mediante post-entrenamiento}
\begin{enumerate}
    \item \textbf{Etapa de preentrenamiento}: se entrenan las \textbf{capacidades visuales fundamentales} del modelo: detección, segmentación, reconocimiento. Estas capacidades están «dormidas»; el modelo ya sabe cómo hacerlo, pero no sabe cuándo hacerlo ni cómo combinarlo.
    \item \textbf{Etapa de post-entrenamiento}: se encapsula una herramienta (como Code Execution) y, mediante aprendizaje por refuerzo, se incentiva al modelo a formar un patrón de comportamiento de \textbf{llamadas a herramientas en múltiples turnos}, con lo que se «libera» la capacidad fundamental.
\end{enumerate}
\end{importantbox}

Este paradigma ofrece dos enseñanzas para la investigación de agentes de IA:

\begin{enumerate}
    \item \textbf{Para quienes diseñan prompts/agentes}: mediante un diseño ingenioso de Visual Prompting y de herramientas, es posible que tareas antes inviables se vuelvan viables sin cambiar los parámetros del modelo
    \item \textbf{Para quienes entrenan modelos}: la calidad de las capacidades fundamentales determina el techo del agente después del post-entrenamiento
\end{enumerate}

\subsection{El papel del Trigger Prompt}

En la demostración oficial de Gemini, el prompt incluye ciertas «palabras disparadoras» para activar la capacidad de Agentic Vision:

\begin{itemize}
    \item "crop out all ..."
    \item "zoom in to see ..."
    \item "annotate on the image ..."
\end{itemize}

\begin{knowledgebox}{El Trigger Prompt es un diseño transitorio}
Si el entrenamiento del modelo fuera suficientemente bueno, en teoría no harían falta estos Trigger Prompt: el modelo debería poder decidir por sí mismo cuándo necesita hacer zoom in y cuándo necesita anotar. La existencia del Trigger Prompt indica que el post-entrenamiento del modelo actual todavía tiene margen de optimización.
\end{knowledgebox}

\subsection{Resumen de la sección}

La enseñanza central de Agentic Vision es que el techo de capacidad de un agente lo determinan las capacidades fundamentales del modelo, y la función del post-entrenamiento es combinar y liberar esas capacidades fundamentales como un comportamiento de razonamiento completo en múltiples turnos.

\section{Análisis de escenarios típicos de aplicación}

\subsection{Escenario uno: conteo preciso (dedos/pedales)}

\textbf{Tarea}: contar cuántos dedos/pedales hay en una imagen.

\textbf{Enfoque de Agentic Vision}:
\begin{enumerate}
    \item El modelo identifica todos los objetivos y genera las coordenadas del BBox de cada uno (capacidad nativa)
    \item Escribe código para dibujar recuadros numerados sobre la imagen, formando un boceto visual
    \item Con base en la imagen anotada, confirma la cantidad
\end{enumerate}

\begin{warningbox}{Discusión sobre lo «tricky»}
El presentador señala que este ejemplo es un tanto «tricky»: al generar los BBox, el modelo ya conoce en realidad la cantidad de objetivos (el número de BBox equivale al número de dedos). La anotación y la segunda confirmación son más bien una «verificación visual» que un verdadero avance en el razonamiento.
\end{warningbox}

\subsection{Escenario dos: análisis de tablas de alta densidad + cálculo}

\textbf{Tarea}: analizar los datos de una tabla en la imagen, realizar un cálculo de normalización y visualizarlo.

\textbf{Problema de los modelos tradicionales}: los VLM presentan con frecuencia alucinaciones (hallucination) en la aritmética visual de varios pasos, sobre todo en operaciones con números de punto flotante.

\textbf{Solución de Agentic Vision}:
\begin{enumerate}
    \item Usa la capacidad nativa de comprensión de gráficas para extraer los datos de la tabla
    \item \textbf{«Descarga» el cálculo a un entorno de Python determinista}: la normalización y demás operaciones matemáticas se completan en el código
    \item Completa la visualización en Matplotlib dentro de Python
\end{enumerate}

\begin{importantbox}{Descarga de cómputo (Computation Offloading)}
Un patrón de diseño importante de Agentic Vision es \textbf{descargar el cálculo preciso de la inferencia del modelo hacia el entorno de ejecución de código}. El modelo se encarga de la comprensión y la planeación, y Python del cálculo preciso. Esto evita de manera eficaz la inestabilidad numérica inherente a los LLM.
\end{importantbox}

\subsection{Escenario tres: clasificación y etiquetado de objetos}

\textbf{Tarea}: clasificar los distintos objetos de la imagen y marcarlos con flechas de diferentes colores hacia sus contenedores correspondientes.

Esta es una tarea integral que exige que el modelo posea al mismo tiempo:
\begin{itemize}
    \item Reconocimiento de objetos (saber qué es)
    \item Razonamiento de clasificación (a qué categoría debe ir)
    \item Anotación visual (dibujar flechas que conecten los objetos con los contenedores)
\end{itemize}

Las pruebas del presentador revelaron que el modelo aún presenta fallas en la lógica de clasificación (por ejemplo, clasifica erróneamente algunos objetos), pero la capacidad general de interacción visual resulta impresionante.

\subsection{Escenario cuatro: gráfica del ciclo de vida animal}

\textbf{Tarea}: recortar de la imagen todos los animales, usar sus fotografías como icon y mostrarlos en una gráfica de Matplotlib ordenados según la duración de su ciclo de vida.

Este caso implica \textbf{un número muy alto de rondas de llamadas a herramientas}: pasos como recortar, ordenar y graficar requieren múltiples ciclos de Think-Act-Observe. La gráfica final resulta a la vez informativa y estéticamente lograda.

\subsection{Resumen de la sección}

Los escenarios de aplicación de Agentic Vision abarcan tareas diversas como el conteo preciso, el análisis de tablas, la clasificación de objetos y la generación de gráficas. Su patrón central es «el modelo se encarga de la comprensión y la planeación, y el código de la ejecución precisa».

\section{Análisis de paquetes capturados: comprensión a nivel de la API}

El expositor realizó un análisis de paquetes capturados de la API de Gemini mediante un hook de solicitudes/respuestas HTTP:

\begin{itemize}
    \item \textbf{Request Body}: contiene la definición de las Tools (Code Execution) y la entrada del usuario (imagen + texto)
    \item \textbf{Response Body}: contiene varias Parts, cada una de las cuales puede ser:
    \begin{itemize}
        \item \texttt{executable\_code}: código Python generado por el modelo
        \item \texttt{code\_execution\_result}: resultado de la ejecución del código (incluidas las imágenes generadas)
        \item \texttt{text}: salida de razonamiento textual del modelo
    \end{itemize}
    \item cada paso del proceso que muestra la interfaz es, en esencia, la visualización de estas Parts
\end{itemize}

\begin{knowledgebox}{Gemini App vs AI Studio}
Al momento de grabarse el video, la Gemini App (la versión para consumidores) aún no ofrecía la Code Execution Tool. Para experimentar Agentic Vision es necesario usar AI Studio (la herramienta para desarrolladores).
\end{knowledgebox}

\section{Implicaciones para la investigación de Agent}

El expositor compartió su reflexión personal sobre las direcciones de investigación de AI Agent:

\begin{enumerate}
    \item \textbf{El valor de un Scaffold Agent puro está disminuyendo}: construir un Agent de propósito general con un desempeño sólido (solid performance) para publicar un artículo tiene cada vez menos valor
    \item \textbf{Enfocarse en un Task/Tool específico}: una dirección de investigación más valiosa es, para una tarea o herramienta concreta, usar el diseño de Agent/Prompting para hacer que \textbf{lo antes inviable se vuelva viable y lo que ya funcionaba mejore de manera considerable}
    \item \textbf{Visual Prompting es una buena dirección}: codifica la información visual en una forma que el Agent puede aprovechar, y es el puente que conecta los modelos de visión con los modelos de lenguaje
    \item \textbf{Agentic RL Training}: establecer las capacidades básicas en el preentrenamiento y, en el Post-Training, usar el aprendizaje por refuerzo para incentivar las llamadas a herramientas de múltiples turnos, liberando así esas capacidades básicas
\end{enumerate}

\section{Síntesis y ampliación}

Gemini Agentic Vision es un trabajo innovador que combina las capacidades visuales fundamentales del modelo con el razonamiento agéntico de múltiples turnos. Su aportación central no radica solo en la funcionalidad misma, sino en revelar una ruta técnica clara:

\begin{enumerate}
    \item \textbf{Capa del modelo}: el preentrenamiento aporta capacidades visuales fundamentales (detección, segmentación, reconocimiento) + Open World Knowledge
    \item \textbf{Capa de herramientas}: se encapsula una Code Execution Tool de propósito general
    \item \textbf{Capa de comportamiento}: mediante Post-Training se incentiva la invocación de herramientas de múltiples turnos al estilo ReAct
    \item \textbf{Capa de aplicación}: el modelo genera de forma autónoma un Visual Scratchpad y razona con base en evidencia visual
\end{enumerate}

\subsection{Lecturas adicionales}

\begin{itemize}
    \item Blog oficial de Gemini Agentic Vision y sus 12 demostraciones de ejemplo
    \item AI Studio: permite probar la Code Execution Tool en línea
    \item El artículo de ReAct (Yao et al., 2023): el fundamento teórico de Reasoning + Acting
\end{itemize}

%% --- Fin del contenido --- %%

\end{document}
